\documentclass[11pt,a4paper]{article}
\usepackage[margin=25mm]{geometry}
\usepackage{amsmath,amssymb,amsthm,hyperref}
\newtheorem{lemma}{Lemma}\newtheorem{theorem}[lemma]{Theorem}
\newcommand{\gen}[1]{\langle#1\rangle}
\title{The binary boundary for full-rank symmetry markings}
\author{Human Compression project: a continuation of the L031 correction}
\date{}
\begin{document}\maketitle
Let $V$ be a vector space over $\mathbf F_2$ with quadratic form $Q$
and polar form $B$. For $Q(x)=1$ the symmetry in $x$ is
\[
 s_x(v)=v+B(v,x)x.
\]
It is an isometry, as substitution into the quadratic identity shows.
It is nonidentity precisely when $x$ is outside the polar radical.
Write $N=\{x\in V:Q(x)=1,\ x\notin\operatorname{rad}B\}$.
We prove the binary nonexistence statement in all dimensions at least
six. Together with the complete two-page companion \emph{Three binary
exceptions to a symmetry exclusion}, this determines exactly the
existence boundary intended in Theorem 1.4 of
Brooksbank--Ferrara--Leemans, arXiv:1709.02219v2.
The companion is required in full for the three positive examples.

\begin{lemma}\label{quad}
Suppose $\dim V=d\geq6$ and $\dim\operatorname{rad}B\leq1$.
Then $N$ spans $V$, and the group generated by all $s_x$, $x\in N$,
acts transitively on $N$.
\end{lemma}
\begin{proof}
The rank $r_B$ of an alternating form is even. Hence $r_B\geq6$
under these assumptions. If $U\leq V$ has codimension $c$, then
\[
 \operatorname{rank}(B|_U)\geq r_B-2c.                 \tag{1}
\]
Indeed, extend a basis of $U$ to one of $V$: adding $c$ rows
and $c$ columns to its Gram matrix increases rank by at most $2c$.

We shall also use the following elementary observation. If $u,v\in U$
satisfy $B(u,v)=1$, then for any $z\in V$,
\[
 Q(z)+Q(z+u)+Q(z+v)+Q(z+u+v)=1.                       \tag{2}
\]
Thus every affine coset of $\gen{u,v}$ contains a vector of
quadratic value one. The brackets here denote linear span.

Suppose $N$ were contained in a hyperplane $U$. By (1),
$B|_U$ has positive rank; choose $u,v\in U$ with $B(u,v)=1$.
For $z\notin U$, partition $z+U$ into its $2^{d-3}$ affine cosets
of $\gen{u,v}$. Equation (2) gives at least $2^{d-3}\geq8$
distinct vectors of quadratic value one outside $U$.
At most one of these can be in the polar radical, since that radical
has dimension at most one. At least one therefore belongs to $N$,
a contradiction. Every proper subspace lies in a hyperplane, so
$N$ spans $V$.

Let $x,y\in N$ be distinct. If $B(x,y)=1$, then $Q(x+y)=1$
and $s_{x+y}(x)=y$. Also $x+y$ is outside the radical, as its
polar product with $x$ is one.
If $B(x,y)=0$, consider the simultaneous equations
\[
                    B(x,z)=B(y,z)=1.                  \tag{3}
\]
The two linear functionals are nonzero. Over $\mathbf F_2$ they
are either independent or equal. Hence (3) is consistent in both
cases, and its solution set is an affine translate of a subspace
$U$ of codimension at most two. By (1), $B|_U$ has rank at least
two. Choose $u,v\in U$ with $B(u,v)=1$. Applying (2) in any
four-point affine coset inside (3) supplies a solution $z$ with
$Q(z)=1$. It is outside the radical since $B(x,z)=1$.
The first case now gives symmetries taking $x$ to $z$ and $z$ to
$y$. Their composition proves transitivity.
\end{proof}

\begin{theorem}\label{no}
Let $d\geq6$ and $\dim\operatorname{rad}B\leq1$.
The full isometry group of $(V,Q)$ cannot be generated by $d$
nonidentity symmetries $s_1,\ldots,s_d$ satisfying
$[s_i,s_j]=1$ for $|i-j|>1$. In particular it has no rank-$d$
string C-group representation by symmetries.
\end{theorem}
\begin{proof}
Suppose such generators exist, and choose $x_i\in N$ with
$s_i=s_{x_i}$. Put $W=\gen{x_1,\ldots,x_d}$. Each generator
preserves $W$ and induces the identity on $V/W$, by the displayed
symmetry formula. Every element of the generated group therefore
induces the identity on $V/W$.

For every $x\in N$, the full isometry group contains $s_x$.
Since $x$ is nonradical, some $v$ satisfies $B(v,x)=1$, and
$(s_x-1)v=x$. Trivial action on $V/W$ thus gives $x\in W$.
Lemma~\ref{quad} implies $W=V$. The $d$ vectors $x_i$ are
consequently a basis.

For distinct basis vectors $x_i,x_j$ of quadratic value one,
$B(x_i,x_j)=1$ prevents their symmetries from commuting:
$s_i s_j(x_i)=x_j$, whereas $s_j s_i(x_i)=x_i+x_j$.
Nonadjacent commutation therefore makes all Gram entries outside
the two adjacent diagonals zero. Form the graph on the basis indices
with an edge $ij$ when $B(x_i,x_j)=1$. It is a subgraph of the
path $1,\ldots,d$.

If that graph were disconnected, the spans of its components would
be a direct sum of nonzero orthogonal subspaces. Every generator
preserves each such subspace, fixing those not containing its root.
Their generated group could not take one basis vector to a basis
vector in another component. But it is the full isometry group,
which acts transitively on $N$ by Lemma~\ref{quad}. This is a
contradiction. The graph is connected, so
\[
 B(x_i,x_j)=1\quad\Longleftrightarrow\quad |i-j|=1.
\]

Put $n=d+1$. In the even-weight space
\[
 E_n=\{z\in\mathbf F_2^n:\textstyle\sum_j z_j=0\},
 \qquad Q_n(z)=|\operatorname{supp}z|/2\pmod2,
\]
the vectors $v_i=e_i+e_{i+1}$ form a basis with precisely these
quadratic values and Gram entries. A quadratic form over $\mathbf F_2$
is determined by its values on a basis and these polar products:
expand $Q(\sum t_i x_i)$ using
$\sum_i t_iQ(x_i)+\sum_{i<j}t_it_jB(x_i,x_j)$.
Thus $x_i\mapsto v_i$ is an isometry $V\longrightarrow E_n$.
In this model $s_i$ is the adjacent coordinate transposition
$(i,i+1)$. Hence the generated group consists exactly of the
coordinate permutations $S_n$.

Since $n\geq7$, choose a vector $u$ with support of size six,
a coordinate $a$ in its support and a coordinate $b$ outside it.
Then $Q_n(u)=1$, while $z=e_a+e_b$ satisfies $B(z,u)=1$.
The isometry $s_u$ sends this weight-two vector to $z+u$, of
weight six. Every coordinate permutation preserves weight.
Therefore $s_u$ is outside the proposed full group $S_n$,
the final contradiction.
\end{proof}

\paragraph{The corrected binary existence statement.}
The companion proves full-isometry symmetry markings of ranks
three, four and five, respectively on $E_4,E_5,E_6$, with the full
intersection property. Their polar radicals have dimension one,
zero and one, exactly the conditions of the supplied theorem.
Theorem~\ref{no} proves that no other dimension $d\geq3$ admits
such a full-rank marking over $\mathbf F_2$. Thus the existential
binary exception set is exactly $\{3,4,5\}$. This does not classify
all individual quadratic forms or all their markings in those
three dimensions.

The source's path argument is retained and made applicable to every
possible full-rank generating set by the spanning and transitivity
lemma. The weight-six symmetry supplies the explicit obstruction
in higher dimension. These are elementary written arguments; no
finite program, external group classification or Lean declaration is
a premise. The positive larger-field theorem and its rank corollaries
remain separate obligations in the full L031 recovery. There is no
whole-original, open-conjecture or priority credit from this note.
\end{document}
